\%============================================================
\chapter{Bài Học: Thành Công Trong Thất Bại, Thất Bại Trong Thành Công (Victory Lives in Defeat, Defeat Hides in Victory)}
\begin{center}
  \textit{\color{truyengold}``Every fall is a lesson; one who never falls has never gone far.''}\\[4pt]
  \textit{(Mỗi lần ngã là một lần học; kẻ không bao giờ ngã là kẻ chưa bao giờ đi xa.)}\\[4pt]
  \small\color{truyenblue}--- Cổ Kim Đông Tây ---
\end{center}
\ngancach

\section{Einstein Trượt Thi Vào Đại Học}
\begin{truyen}{Einstein Trượt Thi Vào Đại Học}{Albert Einstein --- Thụy Sĩ, 1895}
\chuhoa{A}{lbert} Einstein năm 16 tuổi dự thi vào Học viện Bách khoa Zurich và trượt --- kết quả toán vật lý xuất sắc nhưng tiếng Pháp và thực vật học quá kém.\\[4pt]
\textit{(At 16, Albert Einstein sat the entrance exam for ETH Zurich and failed --- outstanding in math and physics, but far too weak in French and botany.)}

Ông phải học thêm một năm ở trường làng Aarau --- điều mà ông sau này gọi là ``may mắn lớn nhất cuộc đời tôi''.\\[4pt]
\textit{(He had to spend an extra year at the village school in Aarau --- something he later called `the greatest stroke of luck in my life.')}

Tại Aarau, ông có thời gian tự do nghĩ về thí nghiệm tưởng tượng: điều gì xảy ra nếu bạn chạy thật nhanh bên cạnh một tia sáng?\\[4pt]
\textit{(At Aarau, he had free time to imagine thought experiments: what would happen if you ran alongside a beam of light?)}

Câu hỏi đó --- nảy sinh từ năm coi là thất bại --- trở thành hạt nhân của Thuyết Tương Đối Đặc Biệt mười năm sau.\\[4pt]
\textit{(That question --- born in what seemed a failed year --- became the nucleus of the Special Theory of Relativity a decade later.)}

Einstein nhận giải Nobel Vật lý năm 1921, nhưng hành trình bắt đầu từ một kỳ thi trượt.\\[4pt]
\textit{(Einstein received the Nobel Prize in Physics in 1921, but the journey started with a failed entrance exam.)}

\end{truyen}
\begin{baihoc}
  Thất bại trong hệ thống hiện có đôi khi là dấu hiệu bạn đang vượt qua hệ thống đó.\\[4pt]
  \textit{(Failing within a system is sometimes a sign that you are outgrowing that system.)}
\end{baihoc}

\section{Michael Jordan Bị Loại Khỏi Đội Tuyển Trường}
\begin{truyen}{Michael Jordan Bị Loại Khỏi Đội Tuyển Trường}{Michael Jordan --- Hoa Kỳ, 1978}
\chuhoa{N}{ăm} lớp 10, Michael Jordan nộp đơn vào đội bóng rổ của trường Emsley A. Laney High School và bị loại vì bị cho là ``quá thấp''.\\[4pt]
\textit{(In 10th grade, Michael Jordan tried out for the varsity basketball team at Emsley A. Laney High School and was cut because he was considered `too short.')}

Ông về nhà khóc một mình trong phòng --- không ai biết điều đó cho đến khi ông tự kể hàng chục năm sau.\\[4pt]
\textit{(He went home and cried alone in his room --- no one knew this until he told the story decades later.)}

Từ đó, ông tập luyện điên cuồng: dậy từ 6 giờ sáng, tập trước khi đến trường, tập sau khi về nhà, tập vào cuối tuần.\\[4pt]
\textit{(From that moment, he trained obsessively: up at 6 a.m., practicing before school, after school, and on weekends.)}

Jordan nói sau này: ``Mỗi khi tôi muốn dừng lại vì mệt mỏi, tôi lại nhớ đến cái bảng danh sách mà tên tôi không có trong đó.''\\[4pt]
\textit{(Jordan later said: `Every time I wanted to stop because I was tired, I thought of that list where my name wasn't on it.')}

Michael Jordan trở thành cầu thủ bóng rổ vĩ đại nhất mọi thời đại --- từ kẻ bị loại trở thành huyền thoại.\\[4pt]
\textit{(Michael Jordan became the greatest basketball player of all time --- from being cut from the team to becoming a legend.)}

\end{truyen}
\begin{baihoc}
  Bị từ chối đúng lúc đúng chỗ có thể là ngọn lửa thắp sáng một sự nghiệp vĩ đại.\\[4pt]
  \textit{(Being rejected at the right time and place can be the fire that ignites a great career.)}
\end{baihoc}

\section{Lần Thứ 1000 Chế Tạo Bóng Đèn Của Edison}
\begin{truyen}{Lần Thứ 1000 Chế Tạo Bóng Đèn Của Edison}{Thomas Edison --- Hoa Kỳ, 1879}
\chuhoa{K}{hi} phóng viên hỏi Edison rằng ông có cảm thấy nản lòng sau hơn 1.000 lần thất bại để tìm ra vật liệu làm sợi đốt bóng đèn không, ông lắc đầu cười.\\[4pt]
\textit{(When a reporter asked Edison if he felt discouraged after more than 1,000 failed attempts to find the right filament for a light bulb, he shook his head and laughed.)}

Edison trả lời: ``Tôi không thất bại 1.000 lần; tôi thành công phát hiện ra 1.000 cách không làm được.''\\[4pt]
\textit{(Edison answered: `I did not fail 1,000 times; I successfully discovered 1,000 ways it doesn't work.')}

Quan điểm này không chỉ là câu nói đẹp --- nó là phương pháp thực sự của ông: mỗi thất bại được ghi chép lại tỉ mỉ và dùng để loại suy cho lần thử tiếp theo.\\[4pt]
\textit{(This attitude was not just a pretty quote --- it was his actual method: each failure was recorded meticulously and used to narrow down the next attempt.)}

Lần thứ 1.001, sợi carbon hóa tre Nhật Bản sáng lên và giữ được hơn 1.200 giờ liên tục.\\[4pt]
\textit{(On the 1,001st try, carbonized Japanese bamboo fiber lit up and held for over 1,200 continuous hours.)}

Mỗi một trong 1.000 lần ``thất bại'' trước đó đều đóng góp một mảnh thông tin không thể thiếu cho thành công cuối cùng.\\[4pt]
\textit{(Each of the 1,000 previous `failures' contributed an indispensable piece of information to the final success.)}

\end{truyen}
\begin{baihoc}
  Thất bại chỉ là dữ liệu --- và dữ liệu chưa bao giờ là kẻ thù của người biết dùng nó.\\[4pt]
  \textit{(Failure is just data --- and data has never been the enemy of those who know how to use it.)}
\end{baihoc}

\section{Alibaba Bị Từ Chối 30 Lần}
\begin{truyen}{Alibaba Bị Từ Chối 30 Lần}{Jack Ma --- Trung Quốc, 1995}
\chuhoa{J}{ack} Ma đã bị Harvard từ chối 10 lần, bị cảnh sát từ chối khi xin làm việc, và bị KFC từ chối khi xin phục vụ bàn cùng 24 ứng viên khác.\\[4pt]
\textit{(Jack Ma was rejected by Harvard 10 times, rejected by the police when applying for work, and rejected by KFC for a waiter position along with 24 other applicants.)}

Năm 1995, công ty đầu tiên của ông thất bại hoàn toàn; ông quay về dạy tiếng Anh với mức lương 12 đô mỗi tháng.\\[4pt]
\textit{(In 1995, his first company failed completely; he returned to teaching English for 12 dollars a month.)}

Năm 1999, ông thành lập Alibaba trong căn hộ nhỏ với 18 người --- mỗi người góp hết tiền tiết kiệm của mình.\\[4pt]
\textit{(In 1999, he founded Alibaba in a small apartment with 18 people --- each contributing all of their savings.)}

Gần chục năm sau đó là những cuộc chiến sinh tử với eBay, với nhà đầu tư tháo lui, với khủng hoảng tài chính toàn cầu.\\[4pt]
\textit{(The following decade was a life-and-death battle against eBay, fleeing investors, and the global financial crisis.)}

Alibaba IPO năm 2014 lập kỷ lục lịch sử với 25 tỷ đô --- và Jack Ma trở thành người giàu nhất Trung Quốc, xuất phát điểm từ tổng cộng hơn 30 lần bị từ chối.\\[4pt]
\textit{(Alibaba's 2014 IPO set a historic record at 25 billion dollars --- and Jack Ma became China's richest man, starting from a total of over 30 rejections.)}

\end{truyen}
\begin{baihoc}
  Số lần bị từ chối không đo lường giới hạn của bạn --- nó đo lường độ kiên trì của bạn.\\[4pt]
  \textit{(The number of times you are rejected does not measure your limits --- it measures your persistence.)}
\end{baihoc}

\section{Nguyễn Bỉnh Khiêm Lui Về Ở Ẩn}
\begin{truyen}{Nguyễn Bỉnh Khiêm Lui Về Ở Ẩn}{Nguyễn Bỉnh Khiêm --- Việt Nam, 1542}
\chuhoa{S}{au} nhiều năm làm quan trong triều Mạc, Nguyễn Bỉnh Khiêm nộp sớ đàn hặc 18 tên gian thần nhưng không được vua nghe, ông xin cáo quan về quê.\\[4pt]
\textit{(After years serving the Mac court, Nguyen Binh Khiem submitted a memorial accusing 18 treacherous officials but the king ignored it, and he resigned and returned home.)}

Nhiều người cho đó là thất bại chính trị, là bỏ cuộc, là kết thúc sự nghiệp.\\[4pt]
\textit{(Many saw this as political failure, as giving up, as the end of a career.)}

Nhưng ông về dựng Am Bạch Vân, mở trường dạy học, viết thơ chữ Nôm và chữ Hán --- và đó mới là đỉnh cao thực sự của cuộc đời ông.\\[4pt]
\textit{(But he returned home, built the White Cloud Hermitage, opened a school, and wrote poetry in both Nom and Chinese script --- and that became the true peak of his life.)}

Từ Am Bạch Vân, ông đào tạo ra những học trò lỗi lạc nhất thế kỷ 16 và để lại kho thơ triết học nhân sinh sâu sắc bậc nhất lịch sử văn học Việt Nam.\\[4pt]
\textit{(From the White Cloud Hermitage, he trained the most distinguished students of the 16th century and left behind a body of philosophical poetry among the deepest in Vietnamese literary history.)}

Ba chúa đều phải đến hỏi ý kiến ông về việc trị quốc --- kẻ ``thất bại'' nơi triều đình trở thành người thầy của các vương triều.\\[4pt]
\textit{(Three lords came to seek his counsel on governance --- the man who `failed' at court became the teacher of dynasties.)}

\end{truyen}
\begin{baihoc}
  Đôi khi rời bỏ hệ thống không phải là thất bại --- đó là sự trưởng thành.\\[4pt]
  \textit{(Sometimes leaving the system is not failure --- it is growth.)}
\end{baihoc}

\section{Abraham Lincoln --- 23 Năm Thất Bại}
\begin{truyen}{Abraham Lincoln --- 23 Năm Thất Bại}{Abraham Lincoln --- Hoa Kỳ, 1832--1860}
\chuhoa{A}{braham} Lincoln thất bại trong kinh doanh năm 1831, thua ứng cử năm 1832, thất bại trong kinh doanh lần nữa năm 1833, và bị thần kinh suy sụp năm 1836.\\[4pt]
\textit{(Abraham Lincoln failed in business in 1831, lost an election in 1832, failed in business again in 1833, and had a nervous breakdown in 1836.)}

Ông thua ứng cử Quốc hội năm 1843, thua thêm lần nữa năm 1848, thua ứng cử Thượng viện năm 1855 và 1858.\\[4pt]
\textit{(He lost congressional bids in 1843 and 1848, lost Senate races in 1855 and 1858.)}

Năm 1860, sau 23 năm liên tiếp thất bại, Lincoln được bầu làm Tổng thống thứ 16 của Hoa Kỳ.\\[4pt]
\textit{(In 1860, after 23 straight years of defeat, Lincoln was elected as the 16th President of the United States.)}

Ông lãnh đạo đất nước qua Nội chiến và ký Tuyên bố Giải phóng Nô lệ --- một trong những hành động pháp lý quan trọng nhất lịch sử nhân loại.\\[4pt]
\textit{(He led the nation through the Civil War and signed the Emancipation Proclamation --- one of the most consequential legal acts in human history.)}

Lincoln từng viết: ``Đường đi của tôi trơn trợt và tôi vấp ngã --- nhưng đó không phải là trượt; tôi đứng dậy và tiếp tục đi.''\\[4pt]
\textit{(Lincoln once wrote: `My path is slippery and I stumbled --- but that was not a slip; I got up and kept going.')}

\end{truyen}
\begin{baihoc}
  Tiểu sử của những người vĩ đại nhất thường là danh sách dài nhất những lần thất bại.\\[4pt]
  \textit{(The biographies of the greatest people are often the longest lists of failures.)}
\end{baihoc}

\section{Phạm Tuân --- Từ Phi Công Bị Thương Đến Nhà Du Hành Vũ Trụ}
\begin{truyen}{Phạm Tuân --- Từ Phi Công Bị Thương Đến Nhà Du Hành Vũ Trụ}{Phạm Tuân --- Việt Nam, 1962--1980}
\chuhoa{P}{hạm} Tuân bắt đầu sự nghiệp phi công năm 1962; năm 1966 máy bay ông bị bắn cháy trên bầu trời Hải Phòng và ông phải nhảy dù thoát thân.\\[4pt]
\textit{(Pham Tuan began his pilot career in 1962; in 1966 his aircraft was shot down over Hai Phong and he parachuted to safety.)}

Thay vì bị chấn thương tâm lý khiến ông bỏ nghề, chính cú rơi đó lại thôi thúc ông học hỏi và hoàn thiện kỹ năng không chiến.\\[4pt]
\textit{(Instead of being traumatized into quitting, that very crash spurred him to study and perfect his aerial combat skills.)}

Năm 1972, ông bắn hạ máy bay B-52 của Mỹ đầu tiên trong lịch sử chiến tranh --- một chiến công làm cả thế giới kinh ngạc.\\[4pt]
\textit{(In 1972, he became the first pilot in history to shoot down a B-52 --- an achievement that astonished the entire world.)}

Và rồi năm 1980, Phạm Tuân trở thành người Việt Nam đầu tiên và người châu Á đầu tiên bay vào vũ trụ.\\[4pt]
\textit{(Then in 1980, Pham Tuan became the first Vietnamese and the first Asian to fly into outer space.)}

Từ chiếc máy bay bị bắn cháy năm 1966 đến vũ trụ năm 1980 --- con đường 14 năm đó chứng minh thất bại chỉ là điểm khởi đầu của một hành trình lớn hơn.\\[4pt]
\textit{(From a burning plane in 1966 to outer space in 1980 --- those 14 years proved that failure is only the starting point of a greater journey.)}

\end{truyen}
\begin{baihoc}
  Cú rơi không đo lường độ cao bạn có thể đạt được --- nó đo lường khả năng bạn đứng dậy.\\[4pt]
  \textit{(A fall does not measure how high you can reach --- it measures your ability to rise.)}
\end{baihoc}

\section{Coco Chanel Xuất Thân Cô Nhi Viện}
\begin{truyen}{Coco Chanel Xuất Thân Cô Nhi Viện}{Coco Chanel --- Pháp, 1895}
\chuhoa{C}{oco} Chanel lớn lên trong cô nhi viện sau khi mẹ mất và cha bỏ đi; bà không có tiền, không có học thức, không có danh phận.\\[4pt]
\textit{(Coco Chanel grew up in an orphanage after her mother died and father left; she had no money, no education, no social standing.)}

Ở thế kỷ đầu 1900, phụ nữ xuất thân như bà hầu như không có cơ hội nào trong xã hội thượng lưu Pháp.\\[4pt]
\textit{(In the early 1900s, a woman of her background had almost no chance in French high society.)}

Nhưng chính vì lớn lên không có gì, bà không bị ràng buộc bởi những quy tắc thời trang của giai cấp thượng lưu.\\[4pt]
\textit{(But precisely because she grew up with nothing, she was not constrained by the fashion rules of the upper class.)}

Trong khi phụ nữ quý tộc mặc áo nịt ngực bó chặt và váy nhiều lớp nặng nề, Chanel cắt may áo đơn giản, thoải mái, thực dụng --- và định nghĩa lại thời trang thế giới.\\[4pt]
\textit{(While aristocratic women wore tight corsets and heavy layered skirts, Chanel cut simple, comfortable, practical clothes --- and redefined world fashion.)}

Thương hiệu Chanel hôm nay là một trong những thương hiệu xa xỉ giá trị nhất thế giới --- bắt đầu từ một đứa trẻ mồ côi không có gì cả.\\[4pt]
\textit{(The Chanel brand today is one of the world's most valuable luxury labels --- begun by an orphan girl who had nothing at all.)}

\end{truyen}
\begin{baihoc}
  Không có gì để mất là sự tự do lớn nhất để đổi mới; nghèo đói đôi khi là cha đẻ của sự sáng tạo.\\[4pt]
  \textit{(Having nothing to lose is the greatest freedom to innovate; poverty is sometimes the parent of creativity.)}
\end{baihoc}

\section{Nguyễn Trãi Sau 10 Năm Tù Túng}
\begin{truyen}{Nguyễn Trãi Sau 10 Năm Tù Túng}{Nguyễn Trãi --- Việt Nam, 1407--1418}
\chuhoa{N}{ăm} 1407, quân Minh xâm chiếm Đại Việt; cha Nguyễn Trãi bị bắt giải về Trung Quốc và ông bị quản thúc ở Đông Quan gần mười năm.\\[4pt]
\textit{(In 1407, Ming forces invaded Dai Viet; Nguyen Trai's father was taken captive to China and Nguyen Trai was placed under house arrest in Dong Quan for nearly ten years.)}

Những năm tù túng đó, ông không lãng phí --- ông quan sát cách cai trị của giặc, nghiên cứu binh thư, suy ngẫm về đường lối cứu nước.\\[4pt]
\textit{(During those years of confinement, he did not waste time --- he observed the enemy's methods of rule, studied military strategy, and pondered ways to save the country.)}

Chính những năm đó đã kết tinh thành bản Bình Ngô Đại Cáo và chiến lược ``tâm công'' --- đánh bằng trí tuệ và nhân nghĩa chứ không chỉ bằng vũ lực.\\[4pt]
\textit{(Those very years crystallized into the Proclamation of Victory over the Wu and the strategy of `winning hearts' --- conquering through wisdom and righteousness, not just force.)}

Khi Lê Lợi dấy binh, Nguyễn Trãi mang theo mười năm trí tuệ tích lũy trong giam cầm làm vũ khí bí mật.\\[4pt]
\textit{(When Le Loi raised his revolt, Nguyen Trai brought with him ten years of wisdom accumulated in captivity as a secret weapon.)}

Bình Ngô Đại Cáo --- được viết từ những tháng năm tăm tối nhất --- trở thành bản tuyên ngôn độc lập vĩ đại nhất lịch sử Việt Nam.\\[4pt]
\textit{(The Proclamation of Victory over the Wu --- written from the darkest years --- became the greatest declaration of independence in Vietnamese history.)}

\end{truyen}
\begin{baihoc}
  Những giai đoạn bị nhốt kín đôi khi giải phóng tinh thần đến độ sâu nhất của nó.\\[4pt]
  \textit{(Periods of confinement sometimes liberate the spirit to its greatest depth.)}
\end{baihoc}

\section{Walt Disney Bị Sa Thải Vì Thiếu Sáng Tạo}
\begin{truyen}{Walt Disney Bị Sa Thải Vì Thiếu Sáng Tạo}{Walt Disney --- Hoa Kỳ, 1919}
\chuhoa{W}{alt} Disney năm 22 tuổi bị sa thải khỏi tờ báo Kansas City Star vì biên tập viên nhận xét ông ``thiếu trí tưởng tượng và không có ý tưởng hay''.\\[4pt]
\textit{(At 22, Walt Disney was fired from the Kansas City Star because an editor said he `lacked imagination and had no good ideas.')}

Công ty phim hoạt hình đầu tiên của ông --- Laugh-O-Gram Studios --- phá sản hoàn toàn năm 1922, ông rơi vào cảnh trắng tay.\\[4pt]
\textit{(His first animation company, Laugh-O-Gram Studios, went completely bankrupt in 1922, leaving him penniless.)}

Ông chuyển đến Hollywood với 40 đô la trong túi và một chiếc vali --- và bắt đầu lại từ đầu.\\[4pt]
\textit{(He moved to Hollywood with 40 dollars in his pocket and one suitcase --- and started over.)}

Mickey Mouse ra đời năm 1928; bộ phim Bạch Tuyết năm 1937 trở thành phim hoạt hình đầu tiên trong lịch sử điện ảnh.\\[4pt]
\textit{(Mickey Mouse was born in 1928; Snow White in 1937 became the first full-length animated film in cinema history.)}

Walt Disney nhận 22 giải Oscar trong sự nghiệp và xây dựng đế chế Walt Disney Company trị giá hàng trăm tỷ đô --- bắt đầu từ kẻ bị cho là ``thiếu trí tưởng tượng''.\\[4pt]
\textit{(Walt Disney won 22 Academy Awards in his career and built the Walt Disney Company, worth hundreds of billions --- starting from someone said to `lack imagination.')}

\end{truyen}
\begin{baihoc}
  Người nói bạn không có tài năng thường là người chưa thấy đủ tầm của tài năng bạn.\\[4pt]
  \textit{(Those who say you lack talent are often those who have not yet seen the full scope of your talent.)}
\end{baihoc}
